\documentclass[11pt]{article}
\usepackage[a4paper,margin=1in]{geometry}
\usepackage{fontspec}
\usepackage[boldfont=false]{xeCJK}
\setCJKmainfont{FandolSong-Regular.otf}
\usepackage{amsmath}
\usepackage{amssymb}
\usepackage{array}
\usepackage{booktabs}
\usepackage{graphicx}
\usepackage{adjustbox}
\usepackage{enumitem}
\setlist[itemize]{topsep=2pt,partopsep=0pt,parsep=0pt,itemsep=2pt,leftmargin=1.4em}
\setlist[enumerate]{topsep=2pt,partopsep=0pt,parsep=0pt,itemsep=2pt,leftmargin=1.6em}
\usepackage{parskip}
\usepackage{fvextra}
\fvset{breaklines=true,breakanywhere=true,fontsize=\small}
\setlength{\parindent}{0pt}

\begin{document}

\begin{center}
  {\LARGE\bfseries An introduction to A Level organic chemistry}\\[0.35em]
  {\large A Level Chemistry}
\end{center}
\vspace{0.6em}

\section*{New functional groups and naming}

At A Level you meet more \textbf{functional group} 官能团 families, including the \textbf{amide} 酰胺 group and \textbf{aromatic} 芳香 compounds (those built on a benzene ring). As before, the functional group decides the properties, and you read it from the general, structural, displayed or skeletal formula.

\subsection*{Naming aromatic compounds}

A \textbf{benzene} 苯 molecule is a ring of six carbons. When you name an aromatic compound, use the \textbf{benzene ring} 苯环 as the parent and number the positions of the substituents. For example, 3-nitrobenzoic acid has a $\text{–NO}_2$ group on carbon 3, and 2,4,6-tribromophenol has three bromine atoms on a phenol ring. You can also name cyclic compounds with a single ring of up to six carbons.

\section*{Two new types of mechanism}

\begin{itemize}
\item \textbf{electrophilic substitution} 亲电取代: an electrophile replaces a hydrogen atom on a benzene ring (this is how benzene reacts).

\item \textbf{addition–elimination} 加成消去: a molecule first adds on, then a small molecule is removed (seen with 2,4-DNPH and with acyl chlorides).

\end{itemize}

\section*{The shape of benzene}

Benzene is a flat, regular hexagon. Each carbon is \textbf{sp²} hybridised, using its three \textbf{hybridisation} 杂化 orbitals to make \textbf{sigma bonds} σ键 to two neighbouring carbons and one hydrogen. This gives the ring of σ bonds.

Each carbon also has one electron left in a p orbital, standing up at right angles to the ring. These p orbitals overlap sideways all the way round, making a single \textbf{delocalised} 离域 \textbf{pi bond} π键 system — a ring of electrons above and below the plane. Because the electrons are shared evenly, all six C–C bonds are the same length, and benzene is very stable.

\section*{Optical isomerism}

Two \textbf{enantiomers} 对映体 (mirror-image isomers) have \textbf{identical} physical and chemical properties, with two exceptions:

\begin{itemize}
\item they rotate \textbf{plane polarised light} 平面偏振光 in opposite directions. A substance that does this is \textbf{optically active} 旋光活性.

\item they may have different effects in living things (biological activity).

\end{itemize}

A \textbf{racemic mixture} 外消旋混合物 is a 50:50 mix of the two enantiomers. It does \textbf{not} rotate plane polarised light, because the two opposite rotations cancel out.

\subsection*{Why chirality matters for drugs}

\textbf{Chirality} 手性 is important when making medicines. A molecule with a \textbf{chiral centre} 手性中心 has two enantiomers, and they can behave very differently in the body — one may cure while the other does harm. So drug makers either:

\begin{itemize}
\item separate a racemic mixture into the two pure enantiomers, or

\item use a \textbf{chiral catalyst} 手性催化剂 to make just the single enantiomer they want.

\end{itemize}


\end{document}
